\section{Extended experimental details}
\label{app:details}

\subsection{Training convergence}

The U-Net with ResNet-34 encoder converges around epoch 60 of 80. We
monitor validation Dice and select the checkpoint with the highest value.
Typical training loss (Dice + BCE hybrid) decreases from $\sim$0.65 at
epoch 1 to $\sim$0.18 at convergence. Validation Dice reaches 0.77 by
epoch 30 and improves by less than 0.01 over the remaining epochs.

Gradient clipping at norm 1.0 is active in roughly 15\% of batches during
the first 10 epochs and less than 2\% after epoch 30, indicating that
training instability is limited to early optimization.

\subsection{Vessel-aware patch sampling}

Class imbalance in retinal vessel segmentation is severe: vessels occupy
10--15\% of pixels in DRIVE images. Uniform random patch sampling would
frequently produce patches with no vessels, wasting training signal. Our
vessel-aware sampling constrains 80\% of patches to contain at least 16
vessel pixels (approximately 0.025\% of a 256$\times$256 patch). The
remaining 20\% are uniformly sampled to maintain exposure to background
regions.

This sampling strategy improves sensitivity by 3--5\% compared to uniform
sampling, with negligible effect on specificity. The 16-pixel minimum
was chosen empirically; lower thresholds (4, 8) produced similar results,
while higher thresholds (32, 64) biased training toward thick vessels and
reduced thin-vessel sensitivity.

\subsection{MC Dropout placement}

We apply Dropout2d (spatial dropout) after the decoder and before the
final 1$\times$1 convolution. Alternative placements were tested:

\begin{itemize}[leftmargin=*]
  \item \textbf{After each encoder block:} Produces higher uncertainty
  magnitude but lower Unc-AUROC (0.72). The dropout perturbations
  propagate through the entire decoder, producing spatially correlated
  noise that does not localize to error-prone regions.

  \item \textbf{After each decoder block:} Similar Unc-AUROC to our
  default (0.76 vs.\ 0.77) but slower inference due to more dropout
  operations per pass.

  \item \textbf{After decoder only (our choice):} Best trade-off
  between Unc-AUROC and inference speed. A single dropout layer
  perturbs the high-resolution feature map just before the
  segmentation head, producing uncertainty that reflects
  representation-level disagreement without the spatial smoothing
  caused by decoder upsampling.
\end{itemize}

\subsection{TTA inverse transform verification}

Geometric TTA requires exact inverse transforms. We verify this by
checking that $\tau_k^{-1}(\tau_k(\mathbf{y})) = \mathbf{y}$ for every
augmentation $\tau_k$ and every test image $\mathbf{y}$. All 6
augmentations (identity, hflip, vflip, rot90, rot180, rot270) are
exact: they permute pixel indices without interpolation, so the
inverse is also an index permutation. This is not the case for
arbitrary rotations or elastic deformations, which require interpolation
and introduce small reconstruction errors.

\section{Additional results}
\label{app:additional_results}

\subsection{Per-image analysis}

\begin{table}[t]
\centering
\caption{Per-image Dice and Unc-AUROC for 5 representative DRIVE test
images, ranked by difficulty (ascending Dice). Hard images (low Dice)
tend to have higher Unc-AUROC, consistent with the cross-dataset finding
that uncertainty quality improves when the model makes more errors.}
\label{tab:per_image}
\small
\setlength{\tabcolsep}{4.6pt}
\resizebox{\textwidth}{!}{%
\begin{tabular}{ccccc}
\toprule
\textbf{Image} & \textbf{Dice} & \textbf{Unc-AUROC (MC)} &
\textbf{Unc-AUROC (TTA)} & \textbf{Uncertainty (mean)} \\
\midrule
Hardest & 0.721 & 0.812 & 0.913 & 0.019 \\
Q1      & 0.748 & 0.789 & 0.895 & 0.015 \\
Median  & 0.765 & 0.770 & 0.878 & 0.012 \\
Q3      & 0.782 & 0.752 & 0.866 & 0.010 \\
Easiest & 0.809 & 0.718 & 0.842 & 0.007 \\
\bottomrule
\end{tabular}%
}
\end{table}

Table~\ref{tab:per_image} shows a clear inverse relationship between
image difficulty and uncertainty quality. On the hardest image (Dice 0.721),
TTA Unc-AUROC reaches 0.913. On the easiest image (Dice 0.809),
Unc-AUROC drops to 0.842. This pattern is beneficial: the images where
deferral matters most are exactly those where the uncertainty signal is
strongest.

Mean uncertainty also correlates with difficulty: 0.019 for the hardest
image vs.\ 0.007 for the easiest. This means mean uncertainty could serve
as a simple image-level difficulty estimator for triage purposes.

\subsection{Failure mode analysis}

We categorize segmentation errors into three types and examine how each
uncertainty method handles them.

\paragraph{Boundary errors.} The most common error type. Vessel boundaries
are annotated at single-pixel precision, but the model's soft predictions
span 2--3 pixels across boundaries. TTA flags these regions with high
variance (boundary pixels often change class under rotation or flip). MC
Dropout assigns moderate uncertainty to boundaries but also to nearby
non-boundary pixels, diluting the signal.

\paragraph{Thin vessel misses.} The model sometimes fails to detect vessels
narrower than 2 pixels, particularly in the periphery. Both MC Dropout and
TTA assign elevated uncertainty to these regions, but TTA does so more
consistently. Under rotation, a thin vertical vessel may align differently
with horizontal convolutional filters, producing disagreement.

\paragraph{False positives on artifacts.} Occasional bright artifacts
(reflections, drusen-like features) are misclassified as vessels. These
are the most dangerous errors for clinical use because the model tends
to be confident (high $\hat{p}$, low $u$). Neither uncertainty method
reliably catches these. Confidence-aware deferral helps slightly because
these pixels, while confident, sometimes have moderate uncertainty, pushing
$s = u \cdot (1-c)$ above threshold. But this failure mode remains the
primary limitation of purely uncertainty-based deferral.

\subsection{Deferral map visualization}

\begin{figure}[t]
\centering
\includegraphics[width=\textwidth]{deferral_map_easy.png}\\[0.5em]
\includegraphics[width=\textwidth]{deferral_map_hard.png}
\caption{Deferral maps comparing global and confidence-aware strategies.
Confidence-aware deferral adapts to image difficulty and concentrates
on pixels near the decision boundary.}
\label{fig:deferral_maps}
\end{figure}

\subsection{Uncertainty histogram analysis}

\begin{figure}[t]
\centering
\includegraphics[width=\textwidth]{qualitative_panel.png}
\caption{Uncertainty distributions for correct vs.\ incorrect pixels.
TTA produces better separation, explaining its higher Unc-AUROC.}
\label{fig:unc_histograms}
\end{figure}

The histograms in Figure~\ref{fig:unc_histograms} make the Unc-AUROC
difference visually intuitive. For MC Dropout, a threshold at any point
on the x-axis will misclassify many correct pixels as uncertain (false
deferrals) or miss many incorrect pixels (missed errors). For TTA, there
exists a threshold range where most incorrect pixels are above and most
correct pixels are below, enabling efficient deferral.

\subsection{Runtime breakdown}

\begin{table}[t]
\centering
\caption{Runtime breakdown per image (average over 20 DRIVE test images).
``Preprocessing'' includes resizing and normalization. ``Aggregation''
includes computing mean, variance, and entropy from multiple passes.}
\label{tab:runtime}
\small
\setlength{\tabcolsep}{4.6pt}
\resizebox{\textwidth}{!}{%
\begin{tabular}{lcccc}
\toprule
\textbf{Component} & \textbf{MC Drop.\ ($T$=30)} & \textbf{TTA ($K$=6)} &
\textbf{Ensemble ($N$=5)} & \textbf{Determ.} \\
\midrule
Preprocessing     & 0.005\,s & 0.005\,s & 0.005\,s & 0.005\,s \\
Forward passes    & 0.780\,s & 0.110\,s & 3.850\,s & 0.065\,s \\
Aggregation       & 0.030\,s & 0.012\,s & 0.040\,s & --- \\
Deferral scoring  & 0.005\,s & 0.003\,s & 0.005\,s & --- \\
\midrule
\textbf{Total}    & \textbf{0.820\,s} & \textbf{0.130\,s} &
\textbf{3.900\,s} & \textbf{0.070\,s} \\
\bottomrule
\end{tabular}%
}
\end{table}

Forward passes dominate runtime for all methods. TTA's advantage comes
from requiring only 6 passes versus 30 for MC Dropout. Aggregation and
deferral scoring are negligible ($<$ 40\,ms).

\subsection{Statistical significance}

\begin{table}[t]
\centering
\caption{Paired $t$-test $p$-values for pairwise method comparisons on
DRIVE test set (20 images). Comparisons with $p < 0.01$ are marked **.
MC Dropout vs.\ Deterministic: uncertainty methods improve all metrics.
MC Dropout vs.\ TTA: TTA wins on Unc-AUROC; segmentation metrics are
comparable.}
\label{tab:significance}
\small
\setlength{\tabcolsep}{4.6pt}
\resizebox{\textwidth}{!}{%
\begin{tabular}{lcc}
\toprule
\textbf{Comparison} & \textbf{Metric} & $p$-\textbf{value} \\
\midrule
MC Drop.\ vs.\ Deterministic & Dice     & $< 0.001$** \\
MC Drop.\ vs.\ Deterministic & AUC      & $< 0.001$** \\
MC Drop.\ vs.\ Deterministic & Unc-AUROC & $< 0.001$** \\
\midrule
MC Drop.\ vs.\ TTA          & Dice      & 0.142 \\
MC Drop.\ vs.\ TTA          & Unc-AUROC & $< 0.001$** \\
\midrule
MC Drop.\ vs.\ Ensemble     & Dice      & $< 0.001$** \\
MC Drop.\ vs.\ Ensemble     & Unc-AUROC & 0.005** \\
\bottomrule
\end{tabular}%
}
\end{table}

The Dice difference between MC Dropout and TTA (0.764 vs.\ 0.768) is
not statistically significant ($p = 0.142$), confirming that TTA's
advantage is in uncertainty quality rather than segmentation quality. The
Unc-AUROC difference is significant ($p < 0.001$), supporting our claim
that TTA produces meaningfully better uncertainty for deferral.

\subsection{Confidence-aware score distribution}

The confidence-aware score $s = u \cdot (1 - c)$ has a different
distribution than raw uncertainty $u$. Raw uncertainty is right-skewed
(most pixels have low uncertainty, a long tail of high uncertainty).
The confidence-aware score is more sharply peaked at zero because many
uncertain pixels are also confident ($c$ close to 1), which suppresses
$s$. This means the confidence-aware threshold $\tau_s$ operates in a
region of the distribution where small changes in $\tau_s$ produce
large changes in the deferral set, making threshold selection more
sensitive. In practice, this sensitivity is manageable because we
optimize $\tau_s$ on a validation set and the optimal value is stable
across cross-validation folds (coefficient of variation $<$ 12\%).

\section{Extended calibration analysis}
\label{app:calibration}

\subsection{Temperature scaling optimization}

Temperature scaling minimizes BCE loss on validation logits using L-BFGS
with a maximum of 50 iterations. Convergence typically occurs within 10
iterations. The optimization landscape is convex (BCE loss is convex in
the logits, and dividing logits by $T$ is a monotonic reparameterization),
so the global optimum is found reliably.

Learned temperatures: $T = 1.19$ (MC Dropout), $T = 1.35$ (TTA). Both
$T > 1$, confirming mild overconfidence. TTA's higher $T$ reflects the
fact that averaging predictions across augmentations tends to sharpen the
mean prediction (the average of several sigmoid outputs is more extreme
than the sigmoid of the average logit), producing slightly more
overconfident outputs than MC Dropout.

\subsection{Bin-level calibration breakdown}

\begin{table}[t]
\centering
\caption{Per-bin calibration analysis for MC Dropout (15 bins). Most
miscalibration concentrates in the lowest and highest confidence bins.
Mid-range bins (0.3--0.7) have near-zero calibration gap.}
\label{tab:bins}
\small
\setlength{\tabcolsep}{4.6pt}
\resizebox{\textwidth}{!}{%
\begin{tabular}{cccccc}
\toprule
\textbf{Bin range} & \textbf{\% pixels} & \textbf{Avg conf} &
\textbf{Avg acc} & \textbf{Gap} & \textbf{Gap (post-cal)} \\
\midrule
$[0.0, 0.07)$  & 82.3\% & 0.013 & 0.016 & 0.003 & 0.004 \\
$[0.07, 0.13)$  & 2.1\%  & 0.095 & 0.112 & 0.017 & 0.015 \\
$[0.13, 0.20)$  & 1.3\%  & 0.163 & 0.186 & 0.023 & 0.019 \\
$[0.20, 0.33)$  & 1.4\%  & 0.261 & 0.278 & 0.017 & 0.013 \\
$[0.33, 0.47)$  & 1.1\%  & 0.398 & 0.405 & 0.007 & 0.005 \\
$[0.47, 0.60)$  & 0.9\%  & 0.531 & 0.522 & 0.009 & 0.008 \\
$[0.60, 0.73)$  & 0.8\%  & 0.665 & 0.648 & 0.017 & 0.014 \\
$[0.73, 0.87)$  & 1.0\%  & 0.796 & 0.782 & 0.014 & 0.011 \\
$[0.87, 0.93)$  & 1.2\%  & 0.898 & 0.891 & 0.007 & 0.006 \\
$[0.93, 1.0]$   & 7.9\%  & 0.978 & 0.984 & 0.006 & 0.005 \\
\bottomrule
\end{tabular}%
}
\end{table}

The dominant contribution to ECE comes from the low-confidence bins
(0.07--0.20), where the model slightly underestimates accuracy. The
high-confidence bin ($> 0.93$) is well calibrated: predicted confidence
0.978 matches observed accuracy 0.984. Temperature scaling reduces the
gap in mid-range bins by 2--4 percentage points but has minimal effect
on the dominant low-confidence and high-confidence bins.

This bin-level analysis explains why overall ECE changes so little:
the bins that contribute most to ECE (low-confidence bins with 2--3\%
of pixels) have gaps of $\sim$0.02, and temperature scaling closes only
a fraction of this gap. The bins with most pixels (near 0 and near 1)
are already well calibrated.

\section{Comparison with deterministic baseline}
\label{app:deterministic}

The deterministic baseline (single forward pass, no uncertainty estimation)
achieves Dice 0.736, AUC 0.944, and ECE 0.245. The high ECE reflects the
fact that a single sigmoid output, without the smoothing effect of
averaging multiple passes, tends to produce extreme probabilities (near
0 or near 1) that are poorly calibrated.

MC Dropout improves Dice by 2.8 percentage points (0.764 vs.\ 0.736), AUC
by 2.1 percentage points, and reduces ECE from 0.245 to 0.038. TTA produces
similar improvements. The segmentation improvement from multiple passes is
well documented \citep{gal2016dropout} and is not the focus of this paper;
we note it for completeness.

The deterministic baseline has Unc-AUROC of exactly 0.500 (random), because
without multiple passes there is no meaningful uncertainty estimate. This
confirms that uncertainty estimation requires stochasticity at test time,
whether from dropout masks (MC Dropout), input transforms (TTA), or separate
models (ensembles).

\section{Notation reference}
\label{app:notation}

\begin{table}[H]
\centering
\caption{Summary of notation used throughout the paper.}
\label{tab:notation}
\small
\setlength{\tabcolsep}{4.6pt}
\resizebox{\textwidth}{!}{%
\begin{tabular}{ll}
\toprule
\textbf{Symbol} & \textbf{Meaning} \\
\midrule
$\mathbf{x}$ & Input image, $\mathbf{x} \in \mathbb{R}^{H \times W \times C}$ \\
$\mathbf{y}$ & Ground truth binary mask, $\mathbf{y} \in \{0,1\}^{H \times W}$ \\
$\hat{\mathbf{p}}$ & Predicted probability map, $\hat{\mathbf{p}} \in [0,1]^{H \times W}$ \\
$\hat{y}_{ij}$ & Hard prediction at pixel $(i,j)$: $\mathbb{1}[\hat{p}_{ij} > 0.5]$ \\
$\mathbf{u}$ & Uncertainty map, $\mathbf{u} \in \mathbb{R}_{\geq 0}^{H \times W}$ \\
$T$ & Number of MC Dropout forward passes \\
$K$ & Number of TTA augmentations \\
$\theta_t$ & Network weights under dropout mask $t$ \\
$\tau_k$ & $k$-th geometric augmentation \\
$\bar{p}_{ij}$ & Mean prediction at pixel $(i,j)$ across passes \\
$\text{MI}_{ij}$ & Mutual information at pixel $(i,j)$ \\
$H(p)$ & Binary entropy: $-p\log p - (1-p)\log(1-p)$ \\
$\delta_{ij}$ & Deferral decision: 1 = accept, 0 = defer \\
$\tau$ & Deferral threshold (global) \\
$\tau_n$ & Deferral threshold for image $n$ (adaptive) \\
$c_{ij}$ & Prediction confidence: $2|\hat{p}_{ij} - 0.5|$ \\
$s_{ij}$ & Confidence-aware deferral score: $u_{ij} \cdot (1 - c_{ij})$ \\
$T_{\text{cal}}$ & Temperature scaling parameter \\
\bottomrule
\end{tabular}%
}
\end{table}
